\section{The critical rational component of every abelian character}
\label{sec:critical-abelian-component-bound}

The following estimate applies to the original integral coefficients of
(FC.1)--(FC.2).  It bounds one rational parity of every character with
trivial cubic factor, including the characters involving the growing
$q$- and $d$-quadratic factors.  The complementary rational parity and
both nontrivial cubic factors remain in a nonnegative arithmetic sum.
The proof does not split a positive derived index into two indices.

Put
\[
 \begin{split}
 F_0&=\mathbb Q(\zeta_{2^{h+2}},\zeta_7,
                       \sqrt{-q},\sqrt{-d}),\qquad
 K=\mathbb Q(\omega),\qquad F=F_0K,\\
 G&=\operatorname{Gal}(F/\mathbb Q),\qquad
 D=[F:\mathbb Q]=96\cdot2^h,\qquad
 \Delta=|\operatorname{disc}F|,\\
 f&=84\cdot2^hqd,\qquad R=2\cdot3\cdot7\cdot269\cdot qd .
 \end{split}                                                    \tag{CP.1}
\]
Here $h\geq1$ and $q,d$ are the actual distinct source primes.
The integer $f$ is a cyclotomic modulus for $F$, and $R$ contains
all fixed finite localization primes.  Let $c\in G$ be complex
conjugation, and put $F^+=F^{\langle c\rangle}$.
For each $\ell$ use exactly the completion $O_\ell$ in (FC.1), and
write $b_\ell=\operatorname{rank}_{\mathbb Z_\ell}O_\ell$.

For a character $\psi$ of $G$ define $\mathcal S^{\mathbb Q}_{\ell,\psi}(e)$
with coefficient
$(O_\ell\otimes\mathbb Q_\ell/\mathbb Z_\ell)(\psi)(-e)$:
take global cohomology unramified outside $\ell R\infty$, impose
zero restriction to the full decomposition group at $\ell$, to
inertia at every other finite localization prime, and to
$H^1(\mathbb R,-)$ at infinity.  At odd $\ell$ the last condition
is automatic.  At $2$ it is part of this definition.

Every character $\eta$ of $\operatorname{Gal}(F/K)$ has exactly two
extensions to $G$, differing by $\chi_{-3}$.  For $n=e+1$ denote
by $\psi_{\eta,e}^{\rm cr}$ the unique one with
$\psi_{\eta,e}^{\rm cr}(c)=(-1)^n$, and by
$\psi_{\eta,e}^{\rm nc}$ the other one.
These are characters in the same original $O_\ell$.

\begin{theorem}[AC-G191: original-integral critical component bound]
\label{AC-G191}
There is an absolute constant $C$ such that, for every actual host
in (CP.1) and every integer $T\geq2$,
\[
 \begin{split}
 \mathcal B_{\rm cr}(F,T)
 &:={\sum_{1\leq e<T}\sum_\ell\frac1{b_\ell}
       \sum_{\substack{\psi\in\widehat G\\\psi(c)=(-1)^{e+1}}}
                      \log\#\mathcal S^{\mathbb Q}_{\ell,\psi}(e)}\\
 &\leq C\mathcal P(F,T),\\
 \mathcal P(F,T)
 &:=DT^2\log(2DTfR)+TD^4\log(2D\Delta R).
 \end{split}                                                    \tag{CP.2}
\]
The sum covers all coefficient primes, including $2$ and $3$,
and every growing finite abelian label of the actual host.

More generally, let $J$ be any subset of the original triples
$(\ell,\eta,e)$ with $1\leq e<T$ in (FC.2), allowing any of the
three cubic factors.  Let $J_0$ be its trivial-cubic triples and
$J_{\rm mix}=J\setminus J_0$.  Define the nonnegative quantity
\[
 \begin{split}
 \mathcal R_J={}&
 \sum_{(\ell,\eta,e)\in J_{\rm mix}}
                  \frac{\log\#\mathcal S_{\ell,\eta}(e)}{b_\ell}\\
 &+\sum_{(\ell,\eta,e)\in J_0}
          \frac{\log\#\mathcal S^{\mathbb Q}_{\ell,
                                  \psi_{\eta,e}^{\rm nc}}(e)}{b_\ell}.
 \end{split}                                                    \tag{CP.3}
\]
For the corresponding original strict sum $\mathcal B_J$,
\[
       |\mathcal B_J-\mathcal R_J|\leq C\mathcal P(F,T).
                                                                  \tag{CP.4}
\]
All sums are finite.  In particular the assertion applies to a
retained set after any of the previously proved pointwise deletions.
No Galois stability of that retained set is required.

For $T=\lfloor u^{399}\rfloor$ and any prescribed source
$\operatorname{WinPair}$ and $\operatorname{PageN}$ data, the
right side of (CP.4) is
\[
       O_\varepsilon\bigl(u^{806+16\varepsilon}(\log u)^2\bigr)
              =o\bigl(\exp(u^\varepsilon/2)\bigr),                 \tag{CP.5}
\]
uniformly in the source pair, page, characters and coefficient
primes.  This is a pointwise estimate at every permitted $u$.
It introduces no separately existential window.
For the positive derived-index sum restricted to the same $J$,
the bound in (CP.5) also holds for $|\mathcal I_J-\mathcal R_J|$,
by the original termwise comparisons specified below.
\end{theorem}

\begin{proof}
\proofstep{1. The actual rational character dictionary.}
The conductor disjointness in (FG.11) gives
$F_0\cap K=\mathbb Q$ and $[F:\mathbb Q]=96\cdot2^h$.
The exponent of $G$ is $3\cdot2^h$, the same exponent used in
(FC.1).  Consequently all its character values already belong
to each specified $O_\ell$, with $b_\ell\leq2^h\leq D$.
Restriction $\widehat G\to\widehat{\operatorname{Gal}(F/K)}$
is two-to-one, with kernel $\{1,\chi_{-3}\}$.
Since $\chi_{-3}(c)=-1$, precisely one lift has each parity.
This proves the dictionary above with the original coefficient
order, including primes dividing the character orders.

\proofstep{2. Finite integral errors over the rational base.}
We use the finite-coefficient kernel comparison (FC.10)--(FC.12),
including its localization square, over $\mathbb Q$.
Its proof is the same long-exact-sequence argument: a finite
kernel $C$ costs global $H^1(C)$ and $H^2(C)$ and local $H^1(C)$.
At the real place the extra term is at most $\log\#C$.
The bound (FG.2) over the imaginary quadratic field $K$ gives
the same bound over $\mathbb Q$ up to an absolute constant.
Indeed the Hochschild--Serre spectral sequence for $K/\mathbb Q$
bounds each rational global $H^i(C)$, $i=1,2$, by products of
$H^j(K,C)$, $j\leq i$, and finite $C_2$-cohomology of those groups.
For a finite $C_2$-module $M$, each $H^j(C_2,M)$ has order at most
$\#M$; also $\#H^0(K,C)\leq\#C$.  The local restriction
sequences give the same assertion, or one may use directly
the finite local Euler-characteristic and inertia bounds from
the proof of (FG.2).  This proves the claimed rational-base
finite estimate, including $\ell=2$ and infinity.

In particular, each Fourier cokernel used below has action
factoring through $F$, is generated by at most $D^2$ abelian
elements, is killed by $D$, and has cost at most
\[
              C D^4\log(2D\Delta R).                              \tag{CP.6}
\]
Only $\ell=2,3$ occur.  To check uniformity, put
$a=v_\ell(D)$ and use the actual trivializing field
$L=F(\mu_{\ell^a})$.  It has degree at most $D^2$;
the tensor-order discriminant argument (FG.10) gives
\[
       \log|\operatorname{disc}L|
          \leq D\log\Delta+C D^2\log(2D).
\]
There are at most a fixed multiple of $[L:\mathbb Q]$ places
above $\ell R$, since $R$ has at most six prime factors.
Substitution in (FG.2), with $a\log\ell\leq\log D$, proves
(CP.6).  The Tate exponent is arbitrary: the same field kills
its action on $C$.  Division by $b_\ell$ can only improve this
coarse estimate.

\proofstep{3. Aggregate the critical parity integrally.}
Fix $e,\ell$, write $O=O_\ell$, and set $s=(-1)^{e+1}$.
Inside the regular lattice use
\[
 L_s=\{x\in O[G]:cx=sx\}
       =\bigoplus_{g\in G/\langle c\rangle}O([g]+s[cg]),
 \qquad
 L'_s=\bigoplus_{\psi(c)=s}O(\psi).                                \tag{CP.7}
\]
The evaluation Fourier map is equivariant and sends a displayed
basis vector to $(2\psi(g))_\psi$.  Orthogonality gives, for its
inverse over the coefficient field,
\[
          a_g=D^{-1}\sum_{\psi(c)=s}y_\psi\psi(g)^{-1}.
                                                                  \tag{CP.8}
\]
Thus $L_s\to L'_s$ is injective and its cokernel is killed by
$D$, with at most $b_\ell D/2\leq D^2$ abelian generators.
This is an actual signed lattice at $2$, not an integral use of
$(1\pm c)/2$.  Tensoring with $\mathbb Q_\ell/\mathbb Z_\ell$
gives the finite-kernel coefficient sequence in the direction
of (FC.10).  Steps 2--3 therefore compare the critical character
sum, divided by $b_\ell$, with the strict group for $L_s(-e)$,
with total all-prime error bounded by (CP.6).

The signed regular lattice is
$\operatorname{Ind}_{G_{F^+}}^{G_{\mathbb Q}}O(\xi_s)$, where
$\xi_+=1$ and $\xi_-$ is the quadratic character of $F/F^+$.
Shapiro and local Mackey identify all the specified strict
local conditions with the conditions over $F^+$.
The resulting group is $b_\ell$ copies, as a
$\mathbb Z_\ell$-module, of the corresponding group with
coefficient $\mathbb Q_\ell/\mathbb Z_\ell(\xi_s)(-e)$.
Thus the original normalization by $b_\ell$ is exact at this
step.  No different, enlarged, or normalized coefficient ring
has been substituted.

\proofstep{4. Strict duality and the finite local conditions.}
The extension $F/F^+$ is unramified at every odd finite prime.
Indeed $i\in F$, whereas an inertia group at an odd prime acts
trivially on $\mathbb Q(i)$; it therefore cannot contain the
nontrivial element $c$ of $\operatorname{Gal}(F/F^+)$.
For $v\nmid\ell$, the invariants of
$A=\mathbb Q_\ell/\mathbb Z_\ell(\xi_s)(-e)$ under inertia are
either zero or a divisible rank-one module.  If $\ell=2$, the
quadratic character is unramified at such $v$.  If $\ell$ is
odd and the quadratic character ramifies, its inertia
invariants vanish.  In the remaining cases Frobenius acts by
a root of unity times $(Nv)^{-e}$, which is not $1$ since
$e>0$.  Multiplication by its difference from $1$ is surjective
on the divisible invariants.  Hence
\[
                    H^1_{\rm unr}(F^+_v,A)=0.                    \tag{CP.9}
\]
So the inertia-zero and full-local-zero conditions agree here.
At every real place, complex conjugation acts on $A$ by
$s(-1)^e=-1$; its $H^1$ is $A/2A=0$, including at $2$.

Let $S$ contain infinity and the places above $\ell R$.
Poitou--Tate duality for the full local kernel, using modified
real terms, now identifies its Pontryagin dual with
$\operatorname{Sha}^2(G_{F^+,S},\mathbb Z_\ell(\xi_s)(n))$.
In particular its order is bounded by
\[
             \#H^2(G_{F^+,S},\mathbb Z_\ell(\xi_s)(n)).          \tag{CP.10}
\]
This uses the compact/discrete form of Sharifi, Appendix A,
Theorem A.0.6 and Remarks A.0.4/A.0.9.  It is a bound by
integral cohomology, not by a rational regulator image.

\proofstep{5. The exact all-prime cohomological inputs.}
For a number field $E$ and $n\geq2$ put
\[
 \mathcal H_n(E)=\prod_\ell H^2_{\rm et}
                   (\mathcal O_E[1/\ell],\mathbb Z_\ell(n)),
 \qquad w_n(E)=\#H^0(E,\mathbb Q/\mathbb Z(n)).                   \tag{CP.11}
\]
The product is finite.  We use the following precise known
special-value formulas, with their cohomological meaning at $2$.
If $E$ is totally real and absolutely abelian and $n$ is even,
\[
                    \#\mathcal H_n(E)
                       =w_n(E)|\zeta_E(1-n)|.                   \tag{CP.12}
\]
For an absolutely abelian CM field $E$, $n\geq3$ odd, and
$E^+$ its real subfield, let
$V_n(E)=\operatorname{coker}(\mathcal H_n(E^+)\to\mathcal H_n(E))$.
Then
\[
              \#V_n(E)\leq
                w_n(E)|L(E/E^+,\xi_-,1-n)|.                     \tag{CP.13}
\]
The first formula follows from Weibel, Chapter VI,
Theorems 8.7--8.8 and 9.12, or the cohomological formula in
Kolster's Appendix A to Rognes--Weibel, Corollary A.3, together
with the odd-prime theorem.  The second is exactly
Kolster--Sands (2008), Sections 5--6, Corollary 6.2, applied
with the cyclic quadratic CM extension: their equality is
$|L|=2^{[E^+:\mathbb Q]}\delta_n\#V_n(E)/w_n(E)$ with
$\delta_n\in\{1,2\}$.  Their Theorem 6.1 imports Kolster's
higher relative class number formula; its $2$-part is known
for absolutely abelian fields.  These deep theorems are retained
inputs, not newly reproved main conjectures or claims about
all noncritical higher class groups.  In particular no
$K$-group formula with its real $2$-correction omitted is used.

\proofstep{6. The remaining integral sign error at $2$.}
For $n$ odd let $M=\mathcal H_n(F)$.  The image of
$\mathcal H_n(F^+)$ is fixed by $c$.  The kernel of
$M^-\to V_n(F)$ is consequently contained in $M[2]$.
Thus
\[
            \#M^-\leq\#V_n(F)\,\#M[2]
                  \leq w_n(F)|L(F/F^+,\xi_-,1-n)|\,\#M[2].       \tag{CP.14}
\]
This elementary inequality replaces an invalid integral
plus/minus decomposition at the prime $2$.
The group $M[2]$ is the $2$-torsion of
$H^2_{\rm et}(\mathcal O_F[1/2],\mathbb Z_2(n))$.
Its order equals that of the quotient by $2$, which injects
into $H^2_{\rm et}(\mathcal O_F[1/2],\mathbb Z/2(n))$.
Over $F$ this finite sheaf is $\mu_2$.  The Kummer exact
sequence and the Brauer invariant sequence give
\[
                  \#M[2]\leq h_F\,2^D
                          \leq 8^D\Delta\,2^D.                  \tag{CP.15}
\]
Here $h_F$ is the ordinary class number; its bound is (FG.5).
There are at most $D$ primes above $2$, and $F$ has no real
places, so the $2$-torsion of the relevant Brauer group has
order at most $2^D$.  Only this elementary ordinary class
bound is used in (CP.15).

We also need descent of the quadratic twist in (CP.10).
Write $M_{\ell,S}=H^2(G_{F,S_F},\mathbb Z_\ell(n))$.
Since $H^0(G_{F,S_F},\mathbb Z_\ell(n))=0$, the low-degree
Hochschild--Serre sequence has an injection of the kernel of
\[
 H^2(G_{F^+,S},\mathbb Z_\ell(\xi_-)(n))
                       \longrightarrow M_{\ell,S}^{-}
\]
into $H^1(C_2,H^1(G_{F,S_F},\mathbb Z_\ell(n))\otimes\xi_-)$.
This kernel is zero for odd $\ell$; at $2$ its order is at most
\[
                         2^{D/2}w_n(F)_{(2)}.                    \tag{CP.16}
\]
Indeed the inner $H^1$ has rank $D/2$ and torsion of order
$w_n(F)_{(2)}$, by the standard integral cohomology description
in Weibel VI.8.2--8.4.  Its rank and torsion are unchanged by
inverting extra primes when $n>1$, by the localization sequence.
For a finitely generated $\mathbb Z_2[C_2]$-module with free
rank $r$ and finite torsion $T$, the cohomology long exact
sequence bounds $\#H^1(C_2,-)$ by $2^r\#T$.
For the torsion-free quotient this follows from the cyclic
cohomology formula and reduction modulo $2$; for $T$ it
follows directly from the same cyclic formula.  This proves
(CP.16) without an unknown index of eigenlattices.

\proofstep{7. Euler terms, higher roots and special-value size.}
The localization sequence for compact integral cohomology is
\[
 0\longrightarrow H^2_{\rm et}(\mathcal O_E[1/\ell],\mathbb Z_\ell(n))
 \longrightarrow H^2(G_{E,S},\mathbb Z_\ell(n))
 \longrightarrow\bigoplus_{v\in S_{\rm fin},\ v\nmid\ell}
                  H^1(k_v,\mathbb Z_\ell(n-1)).                  \tag{CP.17}
\]
The term before the first group is zero since $n>1$.
Each residue group has order the $\ell$-part of
$(Nv)^{n-1}-1$.  Summing its logarithm over all coefficient
primes and the fixed primes above $R$ gives at most
$(n-1)[E:\mathbb Q]\log R$.
Taking a sign subgroup can only decrease the order of the
finite quotient in (CP.17), so this also applies to $M_{\ell,S}^-$.

For any field of degree at most $D$,
\[
                         \log w_n(E)\leq4nD\log(4nD).            \tag{CP.18}
\]
To see this, if $\ell^a\mid w_n(E)$, the image of the
cyclotomic action modulo $\ell^a$ has exponent dividing $n$.
For odd $\ell$ that image is cyclic, so
$\varphi(\ell^a)\leq n[E:\mathbb Q]$ and $\ell^a\leq2nD$.
For $\ell=2$ the group is a subgroup of
$C_2\times C_{2^{a-2}}$; its order is at most $2n$, giving
$2^a\leq4nD$ (the small exponents satisfy this too).
Thus $w_n(E)$ divides $\lfloor4nD\rfloor!$, proving the bound.

Every Dirichlet character belonging to $F$ is used with its
primitive conductor $f_\psi\mid f$ in its Artin $L$-factor.
The standard Hurwitz--Bernoulli special-value identity gives
\[
          |L(1-n,\psi)|\leq3 f^n n! .                            \tag{CP.19}
\]
For completeness, the Bernoulli recurrence inductively gives
$|B_j|\leq j!$: after division by $j!$, its upper bound is
$\sum_{k=2}^{j+1}1/k!<1$.  The binomial formula then gives
$|B_n(x)|\leq3n!$ for $0\leq x\leq1$.
Substitute in
$L(1-n,\psi)=-f_\psi^{n-1}\sum_{a=1}^{f_\psi}
\psi(a)B_n(a/f_\psi)/n$.
For the trivial primitive character the identity is
$\zeta(1-n)=-B_n/n$, with the same bound.
The exact formulas are DLMF 24.2.5, 24.5.3,
25.11.14 and 25.15.3.  This is an archimedean bound for
the whole algebraic critical product; it does not estimate
an individual noncritical regulator.

The products for $\zeta_{F^+}(1-n)$ with $n$ even and for
$L(F/F^+,\xi_-,1-n)$ with $n$ odd each have $D/2$ factors.
Their factors are nonzero by the primitive functional equation
and the Euler product in the half-plane $n>1$.
Hence each product has logarithmic absolute value at most
\[
                  \frac D2\bigl(\log3+n\log f+\log(n!)\bigr).
                                                                  \tag{CP.20}
\]

\proofstep{8. The all-prime bound for the signed regular group.}
If $n$ is even, then $s=+$ and (CP.10) is untwisted
cohomology over the real field $F^+$.  Apply (CP.12),
(CP.17)--(CP.20).  If $n$ is odd, apply (CP.14)--(CP.17)
and then (CP.18)--(CP.20).  In both cases the sum over all
coefficient primes of the logarithms of the signed regular
strict group, normalized by its original coefficient rank,
is at most
\[
                 C\{nD\log(2nDfR)+\log\Delta\}.                 \tag{CP.21}
\]
The bound for $h_F$ in (CP.15) and the factor $2^{D/2}$ in
(CP.16) are included.  The exact all-prime special-value
formulas justify this summation; no unproved equality of
one chosen character's $\ell$-adic lattice with an algebraic
special value has been used.  Adding (CP.6), and summing
$n=e+1$ for $1\leq e<T$, proves (CP.2).

\proofstep{9. Compare the two rational lifts in the original lattice.}
For an original trivial-cubic $\eta$, induction from $K$
gives the rank-two integral lattice
$\operatorname{Ind}_{G_K}^{G_{\mathbb Q}}O_\ell(\eta)$.
Its evaluation map to the direct sum of the two extensions
$O_\ell(\psi_{\eta,e}^{\rm cr})\oplus
 O_\ell(\psi_{\eta,e}^{\rm nc})$ has determinant a unit
times $2$.  It is an isomorphism for odd $\ell$; for $2$
its cokernel is killed by $2$ and has at most $b_2$ abelian
generators.  It is trivialized by $F$, including every Tate
twist modulo $2$.

Shapiro and Mackey identify the induced global and all local
conditions with the original $K$-group (FC.2), also at the
real place, which induces from a complex place.  The same
finite-kernel argument therefore gives the exact equality
for odd $\ell$ and a bounded error at $2$ between
\[
 \frac{\log\#\mathcal S_{\ell,\eta}(e)}{b_\ell}
 \quad\hbox{and}\quad
 \frac{\log\#\mathcal S^{\mathbb Q}_{\ell,
                  \psi_{\eta,e}^{\rm cr}}(e)
       +\log\#\mathcal S^{\mathbb Q}_{\ell,
                  \psi_{\eta,e}^{\rm nc}}(e)}{b_\ell}.            \tag{CP.22}
\]
Summed over any subset of the $D/2$ characters and the given
weights, that finite error is at most
$CTD^4\log(2D\Delta R)$, by Step 2; the rank bound here is
smaller than in (CP.7).  This proves finiteness of the rational
strict groups as well, using the known finiteness in (FC.2).
The all-prime support is finite by that same comparison.
For any retained set $J$, its critical terms form a subset of
the nonnegative full sum in (CP.2).  Equation (CP.22) and
(CP.2) now prove (CP.4), with no stability condition on $J$.

\proofstep{10. The actual growing data and the index boundary.}
The established source estimates give
$D=O_\varepsilon(u^{2+8\varepsilon})$ and
$\log\Delta=O_\varepsilon(u^{2+41\varepsilon/5}(\log u)^2)$.
Also $\log f+\log R\ll\log(2\Delta)$, since the fixed
cyclotomic and quadratic conductors divide the discriminant
up to harmless fixed factors.  The two terms in (CP.2) are
therefore bounded, with $T\leq u^{399}$, by
\[
 O_\varepsilon\left(
  \{u^{802+81\varepsilon/5}+u^{409+201\varepsilon/5}\}
                                      (\log u)^2\right).
                                                                  \tag{CP.23}
\]
Both exponents are smaller than $806+16\varepsilon$ in the
source range $0<\varepsilon<1/52$.  This proves (CP.5), with
the same uniformity as (FG.11) and its displayed source-window
inequalities in the proof of AC-K19.

Finally, for a retained set $J$ define its positive derived
index sum exactly as in (TC.1), restricting the triple sum to
$J$.  The termwise local comparison of AC-K26-local and the
nonnegative torsion bound (EB.3) imply
\[
       |\mathcal I_J-\mathcal B_J|
          \leq O_\varepsilon
                  (u^{806+16\varepsilon}(\log u)^2).             \tag{CP.24}
\]
Indeed the absolute local errors for a subset are bounded by
the full error sum, and
$0\leq\log\#H^2-\log\#k_O\max(0,i)\leq
\log\#H^1_{\rm tors}$ holds term by term.
Combining (CP.24) with (CP.4) gives a bounded comparison with
the explicit nonnegative $\mathcal R_J$.  This does not
identify either rational summand with a separate positive
derived index.  Any application to the current omitted set
must retain its original omitted-cost accounting, its single
finite gluing cost and its same $W$; the present pointwise
polynomial cost can be absorbed within any fixed positive
relative slack on that $W$.
\end{proof}

\paragraph{Remaining arithmetic and exact source scope.}
The unbounded part of (CP.3) still contains every retained
nontrivial cubic character and the noncritical rational lift
of every retained abelian character.  No bound for that
remaining sum, no AM-H04.4 conclusion and no mixed realization
statement is asserted here.  The geometric dictionary is the
finite arithmetic cover and its signed integral induced
lattice; it adds no surface, Hodge class, or numerical-square
construction.  The proof's new estimates and lattice arguments
are complete above.  The deep integral duality and critical
special-value theorems are imported at the exact stated scope.
The source receipt distinguishes the complete printed applied
proofs read from their deeper referenced theorems, and records
that the attempted direct ICTP download returned HTTP 404.
